\chapter{Discuss{\~a}o Cr{\'i}tica}
\label{cap6}
\justifying

{\sloppy Este cap{\'i}tulo analisa de forma cr{\'i}tica os resultados apresentados no Cap{\'i}tulo 5, confrontando-os com a literatura (Cap{\'i}tulo 2) e avaliando a robustez da solu{\c{c}}{\~o}o VotoChain frente a poss{\'i}veis amea{\c{c}}as, limita{\c{c}}{\~o}es metodol{\'o}gicas e validade externa. A discuss{\~a}o est{\'a} organizada em quatro {\'a}reas principais: (i) limita{\c{c}}{\~o}es t{\'e}cnicas, (ii) amea{\c{c}}as {\`a} validade, (iii) compara{\c{c}}{\~o}es com trabalhos relacionados e (iv) implica{\c{c}}{\~o}es te{\'o}ricas e pr{\'a}ticas.\par}

\section{Limitações Técnicas}
\begin{enumerate}
    \item \textbf{Escalabilidade de rede} – Os experimentos foram realizados com 150 eleitores virtuais. Embora o Mix-Net tenha apresentado custos de gas aceitáveis, a carga de transações para eleições nacionais (milhares a dezenas de milhares de eleitores) pode gerar congestionamento na rede Sepolia e elevar drasticamente a latência. Estudos recentes (e.g., \cite{hajian_berenjestanaki_blockchain-based_2024}) sugerem a necessidade de soluções de camada-2 (Rollups) para mitigar esse efeito.
    \item \textbf{Custo de ZK-SNARKs} – O consumo de ~150 k gas por voto torna o ZK-SNARK invi{\'a}vel para grandes elei{\c{c}}{\~o}es sem otimiza{\c{c}}{\~o}es (plonk, halo) ou subs{\'i}dios de gas. A escolha final do mecanismo depender{\'a} do balan{\c{c}}o entre privacidade percebida e viabilidade econ{\'o}mica.
    \item \textbf{Dependência de provedores externos} – O nó Sepolia foi acessado via Infura. Caso o serviço apresente indisponibilidade, a aplicação inteira fica indisponível. Uma solução de fallback (nós auto-hospedados) ainda não foi implementada.
    \item \textbf{Usabilidade limitada ao desktop} – Apesar da conformidade WCAG 2.1, a UI foi testada apenas em navegadores desktop. A experiência em dispositivos móveis, que são predominantes em contextos eleitorais emergentes, ainda não foi avaliada. Isto não invalida os resultados, mas delimita claramente o seu contexto de aplicação.
\end{enumerate}

\section{Ameaças à Validade}
\subsection{Validade Interna}
\begin{itemize}
    \item **Controlo de vari{\'a}veis** – Todos os experimentos foram executados em um ambiente controlado; varia{\c{c}}{\~o}es de rede (picos de congestionamento) n{\~a}o foram simuladas, o que pode inflar a performance observada.
    \item **Viés de seleção dos participantes** – Os 15 participantes do teste de usabilidade foram recrutados entre colegas de laboratório, possivelmente mais familiarizados com tecnologia, o que pode superestimar a usabilidade.
\end{itemize}

\subsection{Validade Externa}
\begin{itemize}
    \item **Generalização para contextos eleitorais reais** – A simulação não inclui fatores sociopolíticos (pressão externa, coação) que podem impactar a confiança no sistema.
    \item \textbf{Interoperabilidade com sistemas de identidade governamentais} – O modelo de registo de eleitores ainda utiliza endereços Ethereum auto‑gerados; a integração com e‑ID ou documentos oficiais ainda não foi testada.
\end{itemize}

\section{Comparação com Trabalhos Relacionados}
A Tabela~\ref{tab:comparativo_critico} posiciona VotoChain frente às plataformas analisadas no Capítulo 2, destacando as dimensões de anonimato, custo e usabilidade.
\begin{table}[h!]
\centering
\caption{Comparativo cr{\'i}tico entre VotoChain e solu{\c{c}}{\~o}es existentes}
\label{tab:comparativo_critico}
\resizebox{\textwidth}{!}{
\begin{tabular}{lccc}
\toprule
Plataforma & Anonimato & Custo de gas (units) & Usabilidade (SUS) \\
\midrule
Voatz          & Parcial (pseud{\'o}nimo) & 85 000 & 62 \\
Helios         & Criptografia homom{\'o}rfica & 45 000 & 68 \\
POLYAS         & Nenhum (identidade vis{\'i}vel) & 120 000 & 70 \\
VotoChain (Mix‑Net) & Total (Mix‑Net) & 90 450 & 71,4 \\
VotoChain (ZK‑SNARK) & Total (ZK‑SNARK) & 150 200 & 71,4 \\
\bottomrule
\end{tabular}
}
\end{table}

Os resultados confirmam que VotoChain oferece \textbf{anonimato total}, algo ausente nas soluções comerciais, ao custo de gas comparável ao de Voatz. A usabilidade, medida pelo SUS, supera a maioria das plataformas, refletindo o esforço de design WCAG.

\section{Implicações Teóricas}
\begin{itemize}
    \item \textbf{Contribuição ao DSR} – O ciclo iterativo (design → demonstração → avaliação) demonstra a viabilidade de aplicar DSR a sistemas críticos de governança, reforçando a proposta de Hevner et al. de que artefactos tecnológicos podem gerar conhecimento teórico ao revelar trade‑offs entre privacidade e desempenho.
    \item \textbf{Modelo de anonimato híbrido} – A implementação modular de ZK‑SNARK e Mix‑Net fornece um referencial empírico para futuras pesquisas que buscam comparar técnicas de anonimato em blockchains públicas.
    \item \textbf{Adoção de métricas de usabilidade em sistemas de votação} – A inclusão de SUS e WCAG como métricas de avaliação amplia o escopo de avaliação de sistemas eleitorais, tradicionalmente focado apenas em segurança e desempenho.
\end{itemize}

\section{Implicações Práticas}
\begin{itemize}
    \item \textbf{Viabilidade para elei{\c{c}}{\~o}es municipais} – Com custos de gas abaixo de 0,05 USD por voto (Mix‑Net) e lat{\^e}ncia <10 s, VotoChain pode ser adotado em elei{\c{c}}{\~o}es de pequeno a m{\'e}dio porte, onde a transpar{\^e}ncia e an{\'o}nimo s{\~a}o cr{\'i}ticos.
    \item \textbf{Roadmap de integração} – A arquitetura modular permite a substituição de componentes (ex.: camada de anonimato) sem reescrever todo o sistema, facilitando a adaptação a requisitos legais de diferentes jurisdições.
    \item \textbf{Pol{\'i}tica de auditoria cont{\'i}nua} – O uso de eventos blockchain e logs imut{\'a}veis possibilita auditorias em tempo real, atendendo {\`a} demanda de transpar{\^e}ncia citada por reguladores eleitorais e preservando o an{\'o}nimo.
\end{itemize}

\section{Direções Futuras de Pesquisa}
\begin{enumerate}
    \item \textbf{Escalabilidade via Layer-2} – Implementar VotoChain em rollups (e.g., Optimism, zkSync) para reduzir custos de gas e melhorar a latência em eleições de grande escala.
    \item \textbf{Estudos de campo} – Conduzir pilotos reais em organizações não-governamentais ou universidades para validar a aceitação social e a resistência a ataques de coerção.
    \item \textbf{Integração com identidade digital governamental} – Desenvolver protocolos de verificação de identidade (e-ID, DID) que preservem anonimato usando credenciais verificáveis (VCs).
    \item \textbf{Otimizações de ZK-SNARK} – Avaliar novas provas (plonk, halo2) que prometem menor consumo de gas e tempos de geração mais curtos.
    \item \textbf{Acessibilidade m{\'o}vel} – Adaptar a UI para dispositivos m{\'o}veis e conduzir testes de usabilidade com utilizadores de diferentes faixas et{\'a}rias e habilidades.
\end{enumerate}


